\label{chap:pantallas-dualidad}
\index{algebraic screen}
\index{duality residue--quotient}

The integers \(10\), \(11\), \(24\), and \(26\) arise in the construction through
different algebraic procedures: the length of a code, the dimension of a
module, the rank of a lattice, and the dimension of an orthogonal quotient. The
purpose of this section is to reconstruct those procedures and the maps that
relate them. Equality between integers will not suffice to identify the
objects they enumerate.

We shall call an \emph{algebraic screen} a subspace or quotient accompanied
by the map that selects it. This denomination does not introduce a new class of
spaces: it abbreviates the pair formed by the linear object and its projector or
quotient map. With this convention, the reductions

\[
12\longrightarrow11\longrightarrow10,
\qquad
26\longrightarrow24
\]

are statements about defined maps, not numerological decompositions.
The first chain arises from a representation of \(A_5\) marked by the
dodecaphasic vector \(K\); the second, from an isotropic reduction of a
Lorentzian lattice. Both rely on the exceptional branch constructed in
the preceding sections, but they are not the same mechanism.

\section{The perfect quotient of \(11\) coordinates}
\label{sec:golay-once}

Let

\[
C_W=[12,6,6]_3
\]

the extended ternary Golay code. Puncturing a coordinate
produces the ternary Golay code

\[
G_{11}=[11,6,5]_3\subset\mathbb F_3^{11}
\]

\cite{MacWilliamsSloane1977}. The length eleven belongs here to the ambient
space of the code; it is not introduced through a physical interpretation.

\begin{definition}[Ternary ball]
For \(x\in\mathbb F_3^n\) and \(t\geq0\), the Hamming ball of radius \(t\)
centered at \(x\) is
\[
B_t(x)=\{y\in\mathbb F_3^n:d_H(x,y)\leq t\}.
\]
Its cardinality, independent of \(x\), is
\[
V_3(n,t)=\sum_{j=0}^{t}\binom nj2^j.
\]
\end{definition}

\begin{theorem}[Perfection of \(G_{11}\)]
\label{thm:g11-perfecto}
The Hamming balls of radius two centered at the words of \(G_{11}\)
constitute a partition of \(\mathbb F_3^{11}\). In particular,
\[
V_3(11,2)=243,
\qquad
|G_{11}|\,V_3(11,2)=3^{11}.
\]
\end{theorem}

\begin{proof}
The minimum distance of \(G_{11}\) is five. If two balls of radius two,
centered at distinct words \(c,c'\in G_{11}\), had a common
element \(y\), the triangle inequality would give
\[
d_H(c,c')\leq d_H(c,y)+d_H(y,c')\leq4,
\]
in contradiction with \(d_H(c,c')\geq5\). The balls are therefore
disjoint. Their cardinality is
\[
V_3(11,2)
=\binom{11}{0}+2\binom{11}{1}+2^2\binom{11}{2}
=1+22+220=243=3^5.
\]
Since \(G_{11}\) has dimension six, it contains \(3^6\) words. In
consequence,
\[
3^6V_3(11,2)=3^6\cdot3^5=3^{11}
=|\mathbb F_3^{11}|.
\]
The disjoint union of the balls has the cardinality of the ambient space and
therefore coincides with it.
\end{proof}

Perfect decoding assigns to each received word a unique codeword
at distance at most two. Equivalently, the syndrome space has
\(3^{11-6}=243\) elements. This occurrence of \(243\) must be
distinguished from the cubic layer

\[
3\cdot9^2=243
\]

of the completion \(1000=729+243+27+1\). One quantity enumerates syndromes of
a code; the other enumerates positions in a boundary layer. Equality
of cardinalities authorizes a comparison, but does not provide a canonical
bijection between the two realizations.

It is also convenient to separate the metric partition

\[
243=1+22+220
\]

of the typed identity

\[
243=1+22+90+120+10.
\]

In the second expression, \(90=6\binom62\) and
\(120=8\binom62\) count families of incidences in the
hexad--octad extension, whereas \(10\) will be the dimension of a linear screen.
They are not subsets of weights \(0,1,2\) of the same Hamming ball.
The construction of designs underpinning that extension belongs to the
classical theory of Witt systems \cite{AssmusKey1992}.

\section{\(2\) descendants of the extended ternary code}
\label{sec:dos-descendientes-cw}

Length eleven is not the only rank that descends from \(C_W\). The same
code intervenes in the gluing of the root lattice \(A_2^{12}\), whose
rank is twenty-four. After fixing the chamber data and the neighborhood vector
constructed in the exceptional branch, the index-three rootless neighbor is the
Leech lattice \(\Lambda_{24}\). The dependency structure is

\begin{equation}
\label{eq:dos-descendientes-cw}
\begin{gathered}
C_W\xrightarrow{\text{puncturing}}G_{11}\subset\mathbb F_3^{11},\\[2mm]
C_W\xrightarrow{\text{gluing}}N(A_2^{12})
\xrightarrow{\text{neighbor with index }3}\Lambda_{24}.
\end{gathered}
\end{equation}

The first arrow is a code morphism; the second line belongs to
lattice theory. The common ternary source explains the coexistence of
ranks eleven and twenty-four without identifying them. In particular, \(11\) is not
a subdimension of \(24\), nor is \(24\) obtained by completing the
punctured code.

\begin{proposition}[Range bifurcation]
The construction \eqref{eq:dos-descendientes-cw} produces two invariants of
different types:
\[
\operatorname{length}(G_{11})=11,
\qquad
\operatorname{rank}(\Lambda_{24})=24.
\]
Their common provenance is the code \(C_W\); neither equality
implies an isomorphism between the coordinate space of the code and the
real space generated by the lattice.
\end{proposition}

\begin{proof}
Puncturing removes one of the twelve positions and does not alter the code's
dimension six, whence length eleven. Each copy of \(A_2\) has rank
two; therefore, \(A_2^{12}\), its finite-index gluings, and its full-rank
neighbors have rank twenty-four. The objects live, respectively,
in the categories of codes over \(\mathbb F_3\) and integral lattices. An
isomorphism between them would not even be well typed without an additional
functor.
\end{proof}

\section{The icosahedral representation of dimension \(12\)}
\label{sec:representacion-a5-doce}

Reduction \(12\to11\to10\) proceeds from a second realization of the number
twelve. Let \(V_{12}=\mathbb R^{12}\) be the permutation module of \(A_5\) on
the twelve vertices of an icosahedron. Denote this representation by \(\rho\)
and its character by \(\chi_{12}\). To fix unambiguously the correspondence
with the twelve dodecaphasic positions, we realize the vertices as left
cosets \(A_5/H\), where \(A_5\) acts on \(\{0,1,2,3,4\}\) and
\[
 H=\langle(0\ 1\ 2\ 3\ 4)\rangle\cong C_5.
\]
The position \(j\) is identified with the coset \(r_jH\), for the
following representatives, written in image notation:
\[
\begin{array}{c|c@{\qquad}c|c@{\qquad}c|c}
j&r_j&j&r_j&j&r_j\\ \hline
1&(3,2,4,1,0)&5&(3,2,0,4,1)&9&(3,1,2,4,0)\\
2&(3,1,0,2,4)&6&(4,0,1,2,3)&10&(2,1,4,3,0)\\
3&(0,1,3,4,2)&7&(4,3,2,1,0)&11&(3,4,1,2,0)\\
4&(1,4,2,3,0)&8&(0,2,3,1,4)&12&(4,2,1,3,0)
\end{array}
\]
This table is the orientation datum that allows a projector of
\(A_5\) to be applied to the vector \(K\); without it, a coordinate norm would depend on an
unspecified permutation. In the order of conjugacy classes

\[
1A,\quad2A,\quad3A,\quad5A,\quad5B,
\]

the sizes of the classes and the number of fixed points are

\[
\begin{array}{c|rrrrr}
 &1A&2A&3A&5A&5B\\ \hline
|C|&1&15&20&12&12\\
\chi_{12}&12&0&0&2&2.
\end{array}
\]

Indeed, rotations of orders two and three fix no vertices, whereas
each rotation of order five fixes the two vertices of its axis. Write

\[
\phi=\frac{1+\sqrt5}{2},
\qquad
\phi'=\frac{1-\sqrt5}{2},
\qquad
\phi+\phi'=1.
\]

The relevant part of the character table of \(A_5\) is

\[
\begin{array}{c|rrrrr}
 &1A&2A&3A&5A&5B\\ \hline
\chi_1&1&1&1&1&1\\
\chi_3&3&-1&0&\phi&\phi'\\
\chi_{3'}&3&-1&0&\phi'&\phi\\
\chi_4&4&0&1&-1&-1\\
\chi_5&5&1&-1&0&0.
\end{array}
\]

\begin{theorem}[Decomposition of the Icosahedral Module]
\label{thm:descomposicion-a5}
The permutation representation on the twelve vertices decomposes as
\[
V_{12}\cong V_1\oplus V_3\oplus V_{3'}\oplus V_5.
\]
The subspace orthogonal to the uniform vector is
\[
V_{11}:=\mathbf1^\perp
\cong V_3\oplus V_{3'}\oplus V_5
\]
and has dimension eleven.
\end{theorem}

\begin{proof}
The multiplicity of an irreducible character \(\chi\) in \(\chi_{12}\) is
\[
\langle\chi_{12},\chi\rangle
=\frac1{60}\sum_C|C|\chi_{12}(C)\chi(C).
\]
Only \(1A,5A,5B\) contribute. For the characters of dimensions one, three,
three-prime, four, and five one obtains, respectively,
\[
\frac{12+24+24}{60}=1,
\]
\[
\frac{36+24(\phi+\phi')}{60}=1,
\qquad
\frac{36+24(\phi'+\phi)}{60}=1,
\]
\[
\frac{48-24-24}{60}=0,
\qquad
\frac{60}{60}=1.
\]
Therefore, the multiplicities are \((1,1,1,0,1)\). The only vector fixed by
all permutations is the uniform vector
\(\mathbf1=(1,\ldots,1)\). Its orthogonal complement contains the three
nontrivial summands and has dimension \(3+3+5=11\).
\end{proof}

The orthogonal projector onto \(V_{11}\) is explicit:

\[
P_{11}=I_{12}-\frac1{12}J_{12},
\]

where \(J_{12}\) is the matrix all of whose entries are one. Thus, the first
reduction is the map

\[
\mathbb R^{12}\longrightarrow V_{11},
\qquad
x\longmapsto x-\frac1{12}\Bigl(\sum_{j=1}^{12}x_j\Bigr)\mathbf1.
\]

This formula specifies which information is eliminated: exclusively the
uniform component. The equality \(12=1+11\) is a consequence of the
decomposition of the module, not a dimensional convention.

\section[Polarization of the three-dimensional sector]{Polarization of the three-dimensional sector determined by K}
\label{sec:reduccion-once-diez}

Consider the already obtained dodecaphasic vector,

\[
K=(234,543,140,729,659,824,621,58,914,794,146,601),
\qquad
\sum_{j=1}^{12}K_j=6263.
\]

Its zero-sum component is

\[
\widetilde K=P_{11}K
=K-\frac{6263}{12}\mathbf1.
\]

To select one of the two three-dimensional representations, we use the
position--coset bijection fixed in the preceding table. The corresponding central
projector is written

\[
P_{\mathrm{ico},3}=\frac{3}{60}\sum_{g\in A_5}\chi_3(g^{-1})\rho(g).
\]

The sum contains permutation matrices and coefficients in
\(\mathbb Q(\sqrt5)\); consequently, \(P_{\mathrm{ico},3}\) is an exact operator over
that field, satisfies \(P_{\mathrm{ico},3}^2=P_{\mathrm{ico},3}=P_{\mathrm{ico},3}^{\mathsf T}\), and has rank three.

\begin{definition}[Polarizing Direction]
The icosahedral projection of the vector and the line that generates it are
\[
u_K=P_{\mathrm{ico},3}\widetilde K,
\qquad
\ell_K=\mathbb R u_K.
\]
\end{definition}

\begin{lemma}[No polarization degeneracy]
\label{lem:norma-uk}
For the vector \(K\) and the declared position--coset correspondence,
\[
\|u_K\|^2
=\frac{6638585+2275584\sqrt5}{20}>0.
\]
In particular, \(\ell_K\) is a well-defined line.
\end{lemma}

\begin{proof}
Upon substituting the expression for \(P_{\mathrm{ico},3}\) into
\(\|u_K\|^2=\widetilde K^{\mathsf T}P_{\mathrm{ico},3}\widetilde K\), the sum reduces to
\[
\frac3{60}\sum_{g\in A_5}
\chi_3(g^{-1})\,
\widetilde K^{\mathsf T}\rho(g)\widetilde K.
\]
Each scalar product is rational, since \(\rho(g)\) is a
permutation matrix and \(\widetilde K\in\mathbb Q^{12}\). The separation of the two
classes of order five leaves as the only irrational part a multiple of
\(\sqrt5\). The exact sum of the sixty terms is
\[
\frac{6638585}{20}+\frac{2275584}{20}\sqrt5.
\]
Both coefficients are positive and \(\sqrt5>0\); therefore, the norm does not
vanish.
\end{proof}

Exact enumeration of the group and the twelve lateral classes reconstructs
the projector onto \(\mathbb Q(\sqrt5)\), verifies its idempotence, and produces the
preceding norm.

\begin{theorem}[Polarized reduction \(11\to10\)]
\label{thm:once-diez}
Let
\[
V_{10}:=(V_3\cap\ell_K^\perp)\oplus V_{3'}\oplus V_5.
\]
Then
\[
V_{11}=\ell_K\oplus V_{10},
\qquad
\dim V_{10}=10.
\]
The projector onto the ten-dimensional screen is
\[
P_{10}=P_{11}-\frac{u_Ku_K^{\mathsf T}}{\|u_K\|^2}.
\]
\end{theorem}

\begin{proof}
By the \cref{lem:norma-uk}, \(u_K\neq0\) and, by definition, \(u_K\in V_3\).
Hence
\[
V_3=\ell_K\oplus(V_3\cap\ell_K^\perp).
\]
Al anadir the sumandos orthogonal \(V_{3'}\) and \(V_5\), the
\cref{thm:descomposicion-a5} gives
\[
V_{11}=\ell_K\oplus
\bigl((V_3\cap\ell_K^\perp)\oplus V_{3'}\oplus V_5\bigr).
\]
The dimension of the second summand is \(2+3+5=10\). The formula for
\(P_{10}\) subtracts from \(P_{11}\) the rank-one orthogonal projector
onto \(\ell_K\), so its image is precisely \(V_{10}\).
\end{proof}

\begin{corollary}[Exact Decomposition \(1+3+7\)]
\label{cor:descomposicion-1-3-7}
Let
\[
W_7(K):=(V_3\cap\ell_K^\perp)\oplus V_5.
\]
Then
\[
\boxed{
V_{11}=\ell_K\oplus V_{3'}\oplus W_7(K),
\qquad
\dim(\ell_K,V_{3'},W_7)=(1,3,7),}
\]
and
\[
\boxed{
V_{10}=V_{3'}\oplus W_7(K),
\qquad 10=3+7.}
\]
\end{corollary}

\begin{proof}
\(\ell_K\subset V_3\) is a line, so that
\(\dim(V_3\cap\ell_K^\perp)=2\). Therefore,
\(\dim W_7=2+5=7\), and the two decompositions follow from the
\cref{thm:once-diez}.
\end{proof}

\begin{remark}[Condition additional for the reading \(1+3+6\)]
\label{rem:interfaz-1-3-6}
The exact decompositions \(1+3+7\) and \(3+7\) do not by themselves select
a time, an extended three-dimensional space, or a compact sector of rank
six. A physical realization must choose a line
\(\ell_t\subset W_7(K)\), prove
\[
W_7(K)=\ell_t\oplus W_6,
\]
construct a full-rank lattice
\(\Lambda_6\subset W_6\) that is stable under transport, interpret
\(W_6/\Lambda_6\) as the compact sector, and exhibit an intertwiner between the
TRIT temporal reader, \(W_6\), and that lattice. The same transport must
separately preserve the three extended directions of the summand
\(V_{3'}\), without mixing them with \(\ell_t\) or with the compact quotient. Only
under those hypotheses are the decompositions obtained
\[
\boxed{
V_{11}=\ell_K\oplus\ell_t\oplus V_{3'}\oplus W_6,
\qquad 11=1+1+3+6,}
\]
\[
\boxed{
V_{10}=\ell_t\oplus V_{3'}\oplus W_6,
\qquad 10=1+3+6.}
\]
The absence of a natural choice of \(\ell_t\), stability of
\(\Lambda_6\), full rank, separation of the three extended
directions, or the intertwiner would refute this realization, without altering the
algebraic corollary \(1+3+7\).
\end{remark}

The reduction is not \(A_5\)-equivariant in the absence of the marked datum. The
submodule \(V_3\) is irreducible and contains no \(A_5\)-invariant line;
the direction \(\ell_K\) appears after introducing \(K\). The sequence

\[
12\longrightarrow11\longrightarrow10
\]

therefore has two distinct natures: the passage \(12\to11\) is canonical for
the permutation action, whereas \(11\to10\) is a polarization
relative to the vector \(K\) and to the declared position--lateral-class correspondence.

\section{Isotropic reduction of rank \(26\to24\)}
\label{sec:reduccion-26-24}

Let \(U\) be the integral hyperbolic plane with basis \(e_+,e_-\) and Gram matrix

\[
\begin{pmatrix}0&1\\1&0\end{pmatrix}.
\]

It is an even, unimodular lattice of signature \((1,1)\). The lattice

\[
L_{26}:=\Lambda_{24}\oplus U
\]

is even, unimodular, and of signature \((25,1)\); by the classification of even
unimodular indefinite lattices, it is identified with \(II_{25,1}\)
\cite{ConwaySloane1999}.

The adjective \emph{Lorentzian} is used in an algebraic sense: it indicates the
presence of a negative direction in the bilinear form. The terminology
comes from Lorentz and Minkowski geometry
\cite{Lorentz1904,Minkowski1909}; it does not by itself identify the lattice with
a physical spacetime.

\begin{theorem}[Null reduction]
\label{thm:reduccion-nula}
For the primitive isotropic vector \(e_+\in U\), there exists a canonical
isomorphism, once the decomposition \(L_{26}=\Lambda_{24}\oplus U\) has been fixed,
\[
e_+^\perp/\mathbb Ze_+\cong\Lambda_{24}.
\]
\end{theorem}

\begin{proof}
Every element of \(L_{26}\) can be written uniquely as
\(x=\lambda+ae_++be_-\), with
\(\lambda\in\Lambda_{24}\) and \(a,b\in\mathbb Z\). Since the sum is
orthogonal and \(\langle e_-,e_+\rangle=1\),
\[
\langle x,e_+\rangle=b.
\]
Therefore,
\[
e_+^\perp=\Lambda_{24}\oplus\mathbb Ze_+.
\]
The quotient by \(\mathbb Ze_+\) removes the second summand and preserves the
first, thereby defining the asserted isomorphism.
\end{proof}

The dimension decreases by two because the orthogonality condition eliminates
one direction and the quotient eliminates the remaining null direction. This is not
an orthogonal projection onto a positive complement: an
isotropic line is contained in its own orthogonal complement. This is the algebraic reason
why the correct procedure is \(e_+^\perp/\mathbb Ze_+\).

The matrix of the plane \(U\) formally coincides with the matrix that exchanges
the two sectors of bidirectional transport. Adopting it as a Gram form
establishes a concrete lattice realization of that involutor. The reduction
\cref{thm:reduccion-nula} is then exact within that realization;
the matrix coincidence, by itself, does not identify the lattice with a
physical spacetime.

This ambient space is also the one used in the Borcherds construction: the
central-charge twenty-four sector is combined with the rank-two Lorentzian
plane before the reduction leading to the Monster Lie algebra
\cite{Borcherds1992}. The relation relevant here is the explicit lattice
quotient; vertex operator algebra theory follows the classical corridor
described above.


\section{Icosahedral McKay correspondence and type \texorpdfstring{\(E_8\)}{E8}}
\label{sec:tres-ochos}

Three occurrences of the integer eight are involved near this boundary:

\[
8_{\rm oct}=|O_H|,
\qquad
V_8:=V_{3'}\oplus V_5,
\qquad
\operatorname{rank}(E_8)=8.
\]

The first is the cardinality of an octad in the Witt system; the second is
a real module of \(A_5\); the third is the rank of a root system. They are not
instances of a single object. The relevant classical bridge passes through the
binary icosahedral group:

\[
A_5\longleftarrow2A_5\subset SU(2)
\xrightarrow{\text{McKay}}\widetilde E_8
\]

\cite{McKay1980}. In the McKay correspondence, the vertices of the affine
diagram encode irreducible representations of \(2A_5\), and the edges
describe the decomposition of the tensor product with the natural
two-dimensional representation. This construction explains why the icosahedron and
\(E_8\) are related; it does not identify a Witt octad with the module
\(V_8\), or either of them with the lattice \(E_8\).

\begin{proposition}[Categorical separation of three value-eight invariants]
No equality among the three preceding invariants defines an interlinker by itself.
To relate two of them one needs, respectively,
a map from incidence sets to modules, from modules of \(A_5\) to
modules of \(2A_5\), or from representations to lattice data.
\end{proposition}

\begin{proof}
The cardinality of a set, the dimension of a vector space, and the rank
of a lattice are invariants in different categories. An isomorphism in
one of them is not defined for objects in the others. The
McKay correspondence provides one of the intermediate functors, but not
the remaining ones; the coincidence of the integer eight is therefore insufficient.
\end{proof}

\section{Dimensional reduction and the context of M-theory}
\label{sec:interfaz-teoria-m}

The preceding constructions produce two exact occurrences of eleven. The first is the length of the perfect code \(G_{11}\); the second is the dimension of the active module \(V_{11}=\mathbf1^\perp\). These objects are not identified: one is defined over \(\mathbb F_3\) and the other over \(\mathbb R\). Reduction to ten belongs exclusively to the latter and depends on the marked line \(\ell_K\).

Eleven-dimensional supergravity and the network of dualities that led to the
formulation of M-theory provide the physical context in which an
eleventh-dimensional direction can be compactified and produce ten-dimensional
theories \cite{CremmerJuliaScherk1978,HullTownsend1995,Witten1995,Townsend1995}.
The dimensional correspondence suggested by the HMT construction is organized
in the following dictionary:

\[
\begin{array}{>{\raggedright\arraybackslash}p{0.42\textwidth}
                >{\raggedright\arraybackslash}p{0.46\textwidth}}
\toprule
\textbf{Objeto algebraico} & \textbf{Dato requerido en una realización física}\\
\midrule
\(V_{11}=\ell_K\oplus V_{10}\) &
una dirección compacta y una pantalla de diez dimensiones con un mapa de
campos definido\\
\((m,w,R_0)\mapsto(w,m,R_0^{-1})\) &
sectores de momento y enrollamiento, radio físico y escala de cuerda\\
\(\Lambda_{24}\oplus U\cong II_{25,1}\) &
una interpretación dinámica del sector reticular lorentziano\\
\(C_W\to G_{11}\) y \(C_W\to\Lambda_{24}\) &
un entrelazador que preserve incidencia, acción y espectro\\
\bottomrule
\end{array}
\]

In eleven-dimensional supergravity, a complete realization must include,
at least, the metric field \(g_{MN}\), the three-form \(C_{MNP}\), the
gravitino \(\psi_M\), its action, and its supersymmetry transformations. For
the screen \(V_{11}=\ell_K\oplus V_{10}\) to represent a physical
compactification, a map is required that sends those fields to data on
\(\ell_K\) and \(V_{10}\) and that preserves the dynamics and the low-energy
limit. None of those fields enters as a hypothesis in the
\cref{thm:g11-perfecto,thm:once-diez,thm:reduccion-nula,thm:dualidad-visible-memoria}.

Consequently, the interface has two clearly differentiated levels. At the algebraic level, perfection of the code, the decomposition \(12=1+3+3'+5\), the polarization \(11=1+10\), the isotropic reduction \(26\to24\), and the radius involution are proved. At the physical level, these structures provide a domain and compatibility equations for constructing a dictionary with M-theory. The second task adds fields and dynamics; it neither modifies nor is used to prove the first level.

\Needspace{24\baselineskip}
\section{Quotient Diagrams and Rank Reductions}
\label{sec:sintesis-pantallas}

The rank diagram can be summarized without losing its types:

\[
\begin{array}{ccccc}
&&C_W=[12,6,6]_3&&\\[1mm]
&\swarrow&&\searrow&\\[-1mm]
G_{11}\subset\mathbb F_3^{11}&&&&N(A_2^{12})\longrightarrow\Lambda_{24}\\[2mm]
&&V_{12}=\mathbb R^{12}&&\\[-1mm]
&&\downarrow P_{11}&&\\[-1mm]
&&V_{11}=V_3\oplus V_{3'}\oplus V_5&&\\[-1mm]
&&\downarrow P_{10}(K)&&\\[-1mm]
&&V_{10}&&\\[2mm]
&&II_{25,1}=\Lambda_{24}\oplus U&&\\[-1mm]
&&\downarrow\ e_+^\perp/\mathbb Ze_+&&\\[-1mm]
&&\Lambda_{24}.&&
\end{array}
\]

The arrow \(P_{11}\) eliminates the uniform mode; \(P_{10}(K)\) eliminates the
direction selected by \(K\); the null reduction eliminates an
isotropic pair by means of orthogonal complement and quotient. The involution
\(\mathsf T\) acts on an enlarged lattice domain and relates reciprocal
radii. These are the mathematical mechanisms supporting the dimensional
interface. Their explicit formulation makes it possible to distinguish precisely the
rank theorems from the physical interpretations that may be built
upon them.
